\chapter{Teorema fundamental del cálculo}\label{antiderivada}

Este capítulo, que presenta uno de los resultados más importantes del cálculo, comienza con la definición de antiderivadas e integrales indefinidas, que se construye a partir de una tabla de derivadas y una tabla de antiderivadas. Además, se resuelven algunos problemas de valor inicial y se presentan el valor promedio de una función y el teorema del valor medio para integrales; a continuación, se definen las funciones de área de acumulación para motivar el teorema fundamental del cálculo (TFC), parte I; después se trabaja con la regla de derivación que este teorema nos proporciona; luego se muestra el TFC, parte II, para calcular integrales definidas a partir de la tabla de integrales indefinidas.

\section{Antiderivada e integral indefinida}

Sea $f$ una función continua. Una antiderivada de $f$ en un intervalo $I$ es cualquier función derivable $F$ tal que $F^\prime(x)=f(x)$ para todo $x$ en $I$.

\begin{ejemplo}
  \textls[-15]{Una antiderivada de $f(x)=3x^{7}\ $es $F(x)=\frac{3}{8}x^{8}$, pues $F'(x)=(\frac{3}{8}x^{8})'=3x^{7}=f(x)$. Otra antiderivada es $F(x)=\frac{3}{8}x^{8}+3$ y, para cada} constante $C$,\ $F(x)=\frac{3}{8}x^{8}+C$ es también una antiderivada de $f$.
\end{ejemplo}

Propiedades:

\begin{itemize}
  \item Si $F$ es una antiderivada de $f$, entonces $F(x) +C$ \ es también una antiderivada de $f$.

  \item La diferencia entre dos antiderivadas de la misma función es una constante.
\end{itemize}

\begin{defi}
  Dada una función $f$ y $F$ una antiderivada, entonces la familia de todas las antiderivadas de $f$, es decir, $F(x)+C$ con $C$ constante, se llama integral indefinida de $f$ y se denota con
  \[\int f(x)dx\]
\end{defi}

\subsection{Propiedades de la integral indefinida}
Supongamos que $f$ y $g$ son funciones con antiderivadas $F$ y $G$ respectivamente. Entonces
\begin{parrafonumerado}
\item $\displaystyle{\int f^{\prime }(x)dx=f(x)+C}$

\item $\displaystyle{\frac{d}{dx}\left(\int f(x)dx\right)=f(x)}$

\item $\displaystyle{\int \left(f(x)+g(x)\right)dx=\int f(x)dx+\int g(x)dx}$

\item $\displaystyle{\int \alpha f(x)dx=\alpha \int f(x)dx},$ donde $\alpha $ es una constante.
\end{parrafonumerado}

\subsection{Algunas integrales indefinidas}
Las siguientes integrales pueden deducirse usando derivadas de funciones conocidas. Por ejemplo, se sabe que para $p$ una constante diferente de $-1$, $(x^p)^\prime=px^{p-1}$ o también $(x^{p+1})^\prime=(p\,+1)x^{p}$, es decir, se tiene la primera propiedad:

\begin{parrafonumerado}

\item Como $\cfrac{d(x^{p+1})}{dx}= (p+1)x^{p}$, entonces, $(p+1) \displaystyle{\int x^{p}dx}=x^{p+1}$ y, por tanto, $\displaystyle{\int x^{p}dx=\frac{x^{p+1}}{p+1}}$ o, mejor,
  \[
    \int x^{p}dx=\frac{x^{p+1}}{p+1}+C
  \]

\item Como $\cfrac{d\, \sen (kx)}{dx}=k\cos (kx)$, para $k$ una constante diferente de cero, entonces se tiene

  \resizebox{\linewidth}{!}{$
    \int k\cos(kx)\,dx=\operatorname{sen}(kx)\Rightarrow
    k\!\int\!\cos(kx)\,dx=\operatorname{sen}(kx)\Rightarrow
    \int\!\cos(kx)\,dx=\tfrac{1}{k}\operatorname{sen}(kx)
  $}

  o también
  \begin{center}
    $\displaystyle{\int \cos (kx)dx=\frac{1}{k}\sen (kx)+C}$
  \end{center}

\item De manera similar al anterior caso, obtenemos las siguientes fórmulas de integrales indefinidas:
  {\small
    \begin{align*}
      \int \sin (kx)dx&=-\frac{1}{k}\cos (kx)+C\\
      \int \sec ^{2}(kx)dx&=\frac{1}{k}\tan (kx)+C\\
      \int \csc ^{2}(kx)dx&=-\frac{1}{k}\cot (kx)+C\\
      \int \sec (kx) \tan (kx)dx&=\frac{1}{k}\sec (kx)+C\\
      \int \csc (kx) \cot (kx)dx&=-\frac{1}{k}\csc (kx)+C
    \end{align*}
  }
  %$\displaystyle{\int \sin (kx)dx=-\frac{1}{k}\cos (kx)+C}$
  %\end{center}
  %
  %\begin{center}
  %$\displaystyle{\int \sec ^{2}(kx)dx=\frac{1}{k}\tan (kx)+C}$
  %%\end{center}
  %
  %\begin{center}
  %$\displaystyle{\int \csc ^{2}(kx)dx=-\frac{1}{k}\cot (kx)+C}$
  %%\end{center}
  %
  %\begin{center}
  %$\displaystyle{\int \sec (kx) \tan (kx)dx=\frac{1}{k}\sec (kx)+C}$
  %%\end{center}
  %
  %\begin{center}
  %$\displaystyle{\int \csc (kx) \cot (kx)dx=-\frac{1}{k}\csc (kx)+C}$
  %%\end{center}

\item Usando que la derivada de $e^{kx}$ es $k e^{kx}$, se tiene que
  \begin{center}
    $\displaystyle{\int e^{kx}dx=\frac{1}{k}e^{kx}+C}$
  \end{center}

\item Se sabe que {\small $(\ln x)^\prime=1/x$} y que {\small $\ln x$} está definida para {\small $x>0$}; por tanto, {\small $\ln(-x)$} está definida para {\small $x< 0$} y {\small $(\ln(-x))^\prime=\frac{1}{-x}(-1)=\frac{1}{x}$}. Además, {\small $\vert x\vert=x$} si {\small $x>0$} y {\small $\vert x\vert=-x$} si {\small $x<0$}; es decir,
  \begin{center}
    $\displaystyle{\int \frac{1}{x}dx=\ln \left\vert x\right\vert +C}$
  \end{center}

\item Usando las derivadas de las funciones trigonométricas inversas, tenemos las siguientes fórmulas:
  {\small
    \begin{center}
      $\displaystyle{\int \frac{1}{\sqrt{a^{2}-x^{2}}}dx=\sen^{-1} \left(\frac{x}{a}\right)+C}$
    \end{center}

    \begin{center}
      $\displaystyle{\int \frac{1}{x^{2}+a^{2}}dx=\frac{1}{a}\tan^{-1} \left(\frac{x}{a}\right)+C}$
    \end{center}

    \begin{center}
      $\displaystyle{\int \frac{1}{\sqrt{x^{2}-a^{2}}}dx=\frac{1}{a}\sen^{-1} \left(\frac{x}{a}\right)+C}$
    \end{center}
  }
\end{parrafonumerado}

En resumen, se obtiene la lista de integrales indefinidas\label{TablaResumen}:

\vspace{\baselineskip}

\begin{center}\label{TablaIntegrales}
  \begin{NiceTabular}{w{c}{0.8cm}w{c}{8cm}}[cell-space-limits=4pt,hvlines]
    \textbf{N.\textsuperscript{o}} & \textbf{Integral} \\
    1 & $\displaystyle\int m\,dx = mx + C$ \\
    2 & $\displaystyle\int x^{p}\,dx = \frac{1}{p+1}x^{p+1}+C,
    \quad p\neq -1$ \\
    3 & $\displaystyle\int \cos(kx)\,dx = \frac{1}{k}\sin(kx)+C$ \\
    4 & $\displaystyle\int \sin(kx)\,dx = -\frac{1}{k}\cos(kx)+C$ \\
    5 & $\displaystyle\int \sec^{2}(kx)\,dx = \frac{1}{k}\tan(kx)+C$ \\
    6 & $\displaystyle\int \csc^{2}(kx)\,dx = -\frac{1}{k}\cot(kx)+C$ \\
    7 & $\displaystyle\int \sec(kx)\tan(kx)\,dx = \frac{1}{k}\sec(kx)+C$ \\
    8 & $\displaystyle\int \csc(kx)\cot(kx)\,dx = -\frac{1}{k}\csc(kx)+C$ \\
    9 & $\displaystyle\int e^{kx}\,dx = \frac{1}{k}e^{kx}+C$ \\
    10 & $\displaystyle\int \frac{1}{x}\,dx = \ln|x|+C$ \\
    11 & $\displaystyle\int \frac{1}{\sqrt{a^{2}-x^{2}}}\,dx = \arcsin\left(\frac{x}{a}\right)+C$ \\
    12 & $\displaystyle\int \frac{1}{a^{2}+x^{2}}\,dx = \frac{1}{a}\arctan\left(\frac{x}{a}\right)+C$ \\
    13 & $\displaystyle\int \frac{1}{x\sqrt{x^{2}-a^{2}}}\,dx = \frac{1}{a}\arcsec\left(\left|\frac{x}{a}\right|\right)+C$ \\
  \end{NiceTabular}
\end{center}

\vspace{\baselineskip}

\begin{ejemplo}
  Usemos las fórmulas dadas anteriormente para calcular algunas integrales:
\end{ejemplo}

\begin{parrafonumerado}

\item $\int \left(3x^{4}-\frac{8}{x^{4}}+7\sqrt{x}+\frac{8}{\sqrt[3]{x}}\right)dx$

  \vspace{\parskip}

  \emph{Solución}
  \begin{small}
    \begin{align*}
      \int &\left(3x^{4}-\frac{8}{x^{4}}+7\sqrt{x}+\frac{8}{\sqrt[3]{x}}\right)dx \\
      &=3\int x^{4}dx-8\int x^{-4}dx+7\int x^{\frac{1}{2}}dx+8\int x^{-\frac{1}{3}}dx \\
      &=\frac{3}{5}x^{5}-\frac{8}{-3}x^{-3}+\frac{\frac{7}{1}}{\frac{3}{2}}x^{\frac{
      3}{2}}+\frac{8}{\frac{2}{3}}x^{\frac{2}{3}} \\
      &=\frac{3}{5}x^{5}+\frac{8}{3}x^{-3}+\frac{14}{3}x^{\frac{3}{2}}+12x^{\frac{2}{3}}+C
    \end{align*}
  \end{small}

\item $\int \bigl(\sen x+\cos x-3\sen (2x)+4\cos (3x)\bigr)dx$

  \vspace{\parskip}

  \emph{Solución}
  \begin{small}
    \begin{align*}
      \int &\bigl(\sen x+\cos x-3\sen (2x)+4\cos (3x)\bigr)dx\\
      &=\int \sen xdx+\int \cos xdx-3\int \sen (2x)dx+4\int \cos (3x)dx\\
      &=-\cos x+\sen x-3\left(-\frac{1}{2}\cos (2x)\right)+4\left(\frac{1}{3}\sen (3x)\right)\\
      &=-\cos x+\sen x+\frac{3}{2}\cos (2x)+\frac{4}{3}\sen (3x)+C
    \end{align*}
  \end{small}

  \enlargethispage*{\baselineskip}

\item $\int \bigl(\csc ^{2}x+3\sec ^{2}(2x)-8\csc ^{2}(7x)\bigr)dx$

  \vspace{\parskip}

  \emph{Solución}
  \begin{small}
    \begin{align*}
      \int \bigl(\csc ^{2}x+&3\sec ^{2}(2x)-8\csc ^{2}(7x)\bigr)dx\\
      &= \int \csc ^{2}xdx+3\int \sec ^{2}(2x)dx-8\int \csc ^{2}(7x)\\
      &=-\cot x+3\left(\frac{1}{2}\tan (2x)\right)-8\left(-\frac{1}{7}\cot (7x)\right)\\
      &=-\cot x+\frac{3}{2}\tan (2x)+\frac{8}{7}\cot (7x)+C
    \end{align*}
  \end{small}

\item $\int \bigl(\sec (2x)\tan (2x)-5\csc (3x)\cot (3x)\bigr)dx$

  \vspace{\parskip}

  \emph{Solución}
  \begin{gather*}
    \begin{aligned}
      \int \bigl(\sec (2x)\tan (2x)-&5\csc (3x)\cot (3x)\bigr)dx\\
      &=\int \sec (2x)\tan (2x)dx - 5\int \csc (3x)\cot (3x)dx\\
      &=\frac{1}{2}\sec (2x)-5\left(-\frac{1}{3}\csc (3x)\right)+L\\
      &=\frac{1}{2}\sec (2x)+\frac{5}{3}\csc (3x)+L
    \end{aligned}
  \end{gather*}

\item $\int \left(e^{2x}+e^{3x}-e^{4x}-\frac{3}{x}\,\right)dx$

  \vspace{\parskip}

  \emph{Solución}
  \begin{align*}
    \int & \left(e^{2x}+e^{3x}-e^{4x}-\frac{3}{x}\,\right)dx \\
    &=\int e^{2x}dx+\int e^{3x}dx-\int e^{4x}dx-3\int \frac{1}{x}dx\\
    &=\frac{1}{2}e^{2x}+\frac{1}{3}e^{3x}-\frac{1}{4}e^{4x}-3\ln |x|+B
  \end{align*}

\item $\int\left( \frac{4}{\sqrt{16-x^{2}}}+\frac{8}{x\sqrt{x^{2}-9}}-\frac{7}{x^{2}+100}\,\right)dx$

  \vspace{\parskip}

  \emph{Solución}
  \begin{align*}
    \int&\left( \frac{4}{\sqrt{16-x^{2}}}+\frac{8}{x\sqrt{x^{2}-9}}-\frac{7}{%
    x^{2}+100}\,\right)dx\\
    &= 4\int \frac{1}{\underbrace{\sqrt{16-x^{2}}}_{a=4}}dx+8\int \frac{1}{x\underbrace{\sqrt{x^{2}-9}}_{a=3}}dx-7\int
    \frac{1}{\underbrace{x^{2}+100}_{a=10}}dx\\
    &=4\sen^{-1} \frac{x}{4}+\frac{8}{3}\sec^{-1}\frac{x}{3}-\frac{7}{10}%
    \tan^{-1} \frac{x}{10}+C
  \end{align*}
\end{parrafonumerado}

\begin{ejemplo}
  Podemos también usar algunas identidades para solucionar ciertas integrales. Veamos dos ejemplos en particular:
\end{ejemplo}

\begin{parrafonumerado}

\item Calculemos la integral de $\sen ^{2}x$. Inicialmente, manipulemos la identidad de $\cos(2x)$ para ver lo que sucede:
  \begin{align*}
    \cos (2x)&=\cos^{2}x-\sen^{2}x\\
    &=1-\sen^{2}x-\sen^{2}x\\
    &=1-2\sen ^{2}x \\
    \Rightarrow 2\sen^{2}x & = 1-\cos (2x)\\
    \Rightarrow \sen^{2}x & = \frac{1}{2}-\frac{1}{2}\cos (2x)
  \end{align*}

  Con ello,
  \begin{align*}
    \int \sen ^{2}xdx&=\int\left( \frac{1}{2}-\frac{1}{2}\cos(2x)\right)dx \\
    &=\frac{1}{2}x-\frac{1}{2}\times \frac{1}{2}\sen (2x)+C \\
    &=\frac{1}{2}x-\frac{1}{4}\sen (2x)+C
  \end{align*}

\item Ahora calculemos la integral de $\cos ^{2}x$. La identidad que usamos en el ítem anterior nos sirve también en este caso, pero vista de la siguiente manera:
  \begin{gather*}
    \cos (2x)=\cos ^{2}x-\sen ^{2}x
    =\cos^{2}x-(1-\cos ^{2}x)
    =2\cos^{2}x-1 \\
    \Rightarrow \quad \cos(2x)+1=2\cos^{2} \quad
    \Rightarrow \quad \cos^{2}x=\frac{1}{2}\cos(2x)+\frac{1}{2}
  \end{gather*}
  Entonces,
  \begin{align*}
    \int \cos ^{2}xdx&=\int\left( \frac{1}{2}+\frac{1}{2}\cos (2x)\,\right) dx\\
    &=\frac{1}{2}x+\frac{1}{2}\times \frac{1}{2}\sen (2x)+C \\
    &=\frac{1}{2}x+\frac{1}{4}\sen (2x)+C
  \end{align*}

\end{parrafonumerado}

\subsection{Problemas de valor inicial}

Un problema de valor inicial tiene como datos alguna derivada de una función desconocida $\frac{d^{n}y}{dx^{n}}$, algunos valores iniciales de la
función $y$ y algunas de las derivadas $\frac{dy}{dx},\frac{d^{2}y}{dx^{2}},\ldots,\frac{d^{n-1}y}{dx^{n-1}}$ en algún punto $x_{0}.$

Veamos ahora ciertos contextos en los que las integrales juegan un papel importante.

\begin{ejemplo}
  Supongamos que un objeto cae libremente desde una altura inicial $s_{0}$, con una velocidad inicial $v_{0}$ y está sujeto únicamente a la fuerza ejercida por la gravedad, es decir, va en caída libre. Quiere describirse la posición y velocidad del objeto en función del tiempo.
\end{ejemplo}

\emph{Solución.} Veamos de qué manera podemos encontrar las ecuaciones de velocidad y posición del objeto. Supongamos inicialmente que $s(t)$ es la posición del objeto, que $v(t)$ es su velocidad y que $a(t)$ es su aceleración en un tiempo cualquiera $t$. Además, en este contexto se tienen ciertas condiciones iniciales:

\begin{itemize}
  \item Posición inicial: $s(0)=s_{0}$.
  \item Velocidad inicial: $\ v(0)=v_{0}$.
  \item Gravedad: $g=9.8\,\nicefrac{\text{m}}{\text{s}^{2}}$.
  \item Aceleración: $a(t)=-g$.
\end{itemize}

Con $s_0$ y $v_0$ constantes.

El enunciado nos dice que el objeto cae libremente; es decir, la aceleración es $g=9.8\,\nicefrac{\text{m}}{\text{s}^{2}}$, y, dado que la velocidad es la derivada de la posición y la aceleración es la derivada de la velocidad, entonces
\begin{equation}\label{aceleracion}
  a(t)=s''(t)=-g
\end{equation}
o, si se quieren unidades explícitas:
\[a(t)=-9.8\,\nicefrac{\text{cm}}{\text{s}^2}\]

Ahora bien, si queremos la posición del objeto, debemos despejar $s$ de \eqref{aceleracion}, es decir, quitar las derivadas de $s$. Recordemos que eso lo podemos hacer por medio de las integrales: $\displaystyle{\int s^{\prime \prime} (t)dt=s^\prime (t)}$; entonces
\begin{align*}
  s^\prime (t)&=\int s^{\prime \prime} (t)\,dt=\int a(t)dt=\int -9.8dt=-9.8 t+C_1
\end{align*}
es decir, $s^\prime(t)=v(t)=-9.8t+C_1$, pero $v(0)=v_{0}$, luego $v_0=-9.8\times 0+C_1$, entonces $v_0=C_1$; con ello,
\[
  s^\prime(t)=v(t)=-9.8t+v_0
\]
Por otro lado, $\displaystyle{\int s^\prime (t)dt=s(t)}$; entonces
\begin{align*}
  s(t)&=\int s^\prime (t)dt=\int v(t)dt=\int 9.8t+v_{0}dt=-\frac{9.8}{2}t^{2}+v_{0}t+C_1
\end{align*}
pero $s(0)=s_0$, entonces
\[
  s_0=-\frac{9.8}{2}\times 0^{2}+v_{0}\times 0+C_1\, \Rightarrow \, s_0=C_1
\]
y, con ello, la posición de la partícula en las condiciones del enunciado es
\begin{equation}\label{posicion}
  s(t)=-4.9t^{2}+v_{0}t+s_{0}
\end{equation}

\begin{ejemplo}
  \textls[-15]{Encuentre el valor de $f$ tal que $f'''(x)=24x+16\sen (2x)-81\cos (3x)$ sujeta a las condiciones iniciales $f(0)=1$, $f'(0)=2$\; y $f''(0)=0$.}
\end{ejemplo}

\emph{Solución.} Como en el ejemplo anterior, la idea es quitar las derivadas integrando:
\begin{align*}
  f^{\prime \prime }(x)=\int f^{\prime \prime\prime }(x)\, dx &=\int (24x+16\sen (2x)-81\cos (3x))dx\\
  &=12x^{2}-8\cos (2x)-27\sen (3x)+A
\end{align*}
pero $f^{\prime \prime }(0)=0$, entonces $0=-8+A$ y, así, $A=8$. Por otro lado,
\begin{align*}
  f^{\prime }(x)=\int f^{\prime \prime }(x)\, dx &= \int (12x^{2}-8\cos (2x)-27\sen (3x)+8)dx \\
  &=4x^{3}-4\sen (2x)+9\cos (3x)+8x+B
\end{align*}
y se nos dijo que $f^\prime (0)=2$, entonces $2=9+B$ y $B=-7$. Finalmente,
\begin{align*}
  f(x)=\int f^\prime(x) \, dx &=\int (4x^{3}-4\sen (2x)+9\cos (3x)+8x-7)dx \\
  &=x^{4}+2\cos (2x)+3\sen (3x)+4x^2-7x+E
\end{align*}
que, al usar la condición inicial suministrada, $f(0)=1$, se tiene $1=2+E$ y $E=-1$.

Al agrupar toda la información encontrada, tenemos que
\[
  f(x)=x^{4}+2\cos (2x)+3\sen (3x)+4x^2-7x-1
\]

\begin{ejemplo}
  La pendiente de la recta tangente a $y=f(x)$ en el punto $(x,f(x))$ está dada por $f^{\prime }(x)=\cfrac{x^{2}+8}{x^{2}+9}$. Si $f(3)=2$, encuentre $f(x)$.
\end{ejemplo}

\emph{Solución.} Como $f$ es la integral de $f^\prime$ y, además, $\cfrac{x^{2}+8}{x^{2}+9}=1-\cfrac{1}{x^{2}+9}$, entonces
\begin{align*}
  f(x)&=\int f^{\prime }(x)\,dx=\int \frac{x^{2}+8}{x^{2}+9}\, dx \\
  &=\int \left( 1-\frac{1}{x^{2}+9}\right)\,dx \\
  &=x-\frac{1}{3}\tan^{-1} \frac{x}{3}+A
\end{align*}

Usando la condición inicial,
\[
  f(3)=3-\frac{1}{3}\tan^{-1}(1)+A=2
\]
entonces
\begin{align*}
  3-\frac{1}{3}\cdot\frac{\pi }{4}+A&=2\\
  A&=-1+\frac{\pi }{12}
\end{align*}
es decir,
\[
  f(x)=x-\frac{1}{3}\tan^{-1}\left(\frac{x}{3}\right)-1+\frac{\pi}{12}
\]

\begin{ejemplo}
  En $t=1/2$, un tren se aproxima a una estación y comienza a desacelerar llevando una velocidad de $120$ millas por hora. La función de aceleración del tren en millas por hora al cuadrado es $a(t) =-200 t ^{-3}\text{mi}/\text{h}^{2}$, donde $t\geq 1/2.$ ¿Cúal es la distancia recorrida por el tren entre $t=1$ y $t=1.1$ horas?
\end{ejemplo}

\emph{Solución.} Como nos dan la función de aceleración, entonces
\begin{align*}
  v(t) &=\int -200 t^{-3}dt%\\
  %&
  =100t^{-2}+C \
  =\frac{100}{t^{2}}+C
\end{align*}
y como $v\left( 1\right) =120$, entonces $100+C=120$, $C=20$; así, $v(t)=\frac{100}{t ^{2}}+20$.

Ahora integramos la función de velocidad $v(t) $ para obtener la función de posición o distancia:
\begin{align*}
  s(t) =\int \left( 100t^{-2}+20\right)
  dt=-100t^{-1}+20t+B
\end{align*}
La distancia total recorrida por el tren es
\begin{align*}
  s\left( 1.1\right) -s\left( 1\right) &=-100\left( 1.1\right) ^{-1}+20\left(
  1.1\right) +B-\left( -100+20+B\right) \\
  &=11.091\, \text{millas}
\end{align*}

\begin{ejemplo}
  Supongamos que $a(t)=8\sen ^{2}(2t)+16\cos ^{2}(4t)$ es la aceleración de una partícula en movimiento y encontremos su posición si $v(0)=2$ y $s(0)=1$:
\end{ejemplo}

\emph{Solución}
\begin{align*}
  v(t)&=\int a(t)dt=\int \left( 8\left(\frac{1}{2}-\frac{1}{2}\cos(4t)\right)+16\left(\frac{1}{2}+%
  \frac{1}{2}\cos (8t)\right) \right) \,dt\\
  &=\int (12-4\cos (4t)+8\cos (8t))\,dt=12t-\sen (4t)+\sen (8t)+A
\end{align*}
pero $v(0)=2$, entonces $A=2$.

Por otro lado,
\begin{align*}
  s(t)&=\int v(t)dt=\int (12t+2-\sen (4t)+\sen (8t))dt\\
  &=6t^2+2t+\frac{1}{4}\cos(4t)-\frac{1}{8}\cos(8t)+B
\end{align*}
y al usar que $s(0)=1$, se tiene que
\[
  1=\frac{1}{4}-\frac{1}{8}+B \Rightarrow B=1-\frac{1}{8}=\frac{7}{8}
\]
y, con ello, la posición de la partícula es
\[
  s(t)=6t^{2}+2t+\frac{1}{4}\cos (4t)-\frac{1}{8}\cos (8t)+\frac{7}{8}
\]

\section{Teorema del valor medio para integrales}
En la vida cotidiana aparecen preguntas como ¿cuál es la temperatura promedio de un día en una ciudad? o ¿cuál es la velocidad promedio de un auto en un cierto recorrido? En este apartado se define el valor promedio de una función continua en un intervalo cerrado y se presenta y demuestra el teorema del valor medio para integrales, que nos dice que, dada una función continua en un intervalo cerrado, existe al menos un punto del intervalo donde la función evaluada en este punto multiplicada por la longitud del intervalo es igual a la integral de la función en este intervalo.

\subsection{Valor promedio de una función}

Dados $n$ números reales $y_{1},y_{2},\ldots,y_{n}$, el promedio\ $\overline{y}$%
\ de estos números es
\[
  \overline{y}=\frac{y_{1}+y_{2}+\cdots+y_{n}}{n}
\]
Sea $f$ una función continua en el intervalo $[a,b]$. El \textit{promedio de la función} $f$ en el intervalo $[a,b]$ se obtiene al \textit{promediar} las infinitas imágenes $f(x) $ para $x\in \lbrack a,b]$. Hagámoslo siguiendo el camino tomado con las integrales. Consideremos una partición cualquiera del intervalo $[a,b]$: $\left\{ x_{1},x_{2},x_{3},\cdots ,x_{n}\right\} $,
elegimos un punto $x_{k}^{\ast }$ en el $k$-ésimo intervalo $%
[x_{k-1},x_{k}]$ y promediamos las $n$ imágenes:
\begin{equation*}
\left\{ f\left( x_{1}^{\ast
  }\right) ,f\left( x_{2}^{\ast }\right) ,f\left( x_{3}^{\ast }\right) ,\cdots
,f\left( x_{n}^{\ast }\right) \right\}
\end{equation*}
Con ello, el promedio de $f$ en el intervalo $[a,b]$, denotado por $\Pro\left( f\right) $, puede aproximarse por
\begin{align*}
\Pro\left( f\right) &\approx \frac{f\left( x_{1}^{\ast }\right) +f\left(
  x_{2}^{\ast }\right) +f\left( x_{3}^{\ast }\right) +f\left( x_{4}^{\ast
}\right) }{n}\\
&=\frac{\sum_{k=1}^{n}f\left( x_{k}^{\ast }\right) }{n}\\
&=\frac{\frac{b-a}{n}\sum_{k=1}^{n}f\left( x_{k}^{\ast }\right) }{b-a}\\
&=\frac{\sum_{k=1}^{n}f\left( x_{k}^{\ast }\right) \Delta x}{b-a}
\end{align*}
Así, haciendo que $n$ tienda a infinito en esta suma (que es una
suma de Riemann), cuando el límite existe, este es el valor promedio de $f$:
\begin{align*}
\Pro\left( f\right)&=\lim_{n\rightarrow \infty }\frac{\sum_{k=1}^{n}f\left(
x_{k}^{\ast }\right) \Delta x}{b-a}=\frac{1}{b-a}\int_{a}^{b}f\left(
x\right) dx
\end{align*}

\begin{ejemplo}
Sea $T(t) =t^{3}+t$ la temperatura en grados centígrados de una lámina. Halle la temperatura promedio de la lámina para $0\leq t\leq 4$.
\end{ejemplo}

\emph{Solución.}
La temperatura promedio es:
\begin{gather*}
\begin{aligned}
  \Pro(T)&=\frac{1}{4-0}\int_{0}^{4}T(t)dt %\\&
  =\frac{1}{4}\int_{0}^{4}(t^{3}+t)dt=\frac{1}{4}\left( \frac{4^{4}-0^{4}}{4}+\frac{4^{2}-0^{2}}{2}\right) \\
  &=\frac{1}{4}\left(64+8\right)=18\,\degree C
\end{aligned}
\end{gather*}

El teorema del valor medio para integrales nos dice que el valor promedio de una función continua $f$ en un intervalo $[a,b]$ es el valor de la función en algún punto del intervalo $(a,b)$.
\begin{teorema}\label{ValorMedio}
Si $f$ es una función continua en el intervalo $\left[a,b\right]$, para $a$ y $b$ números reales, entonces existe un número real $c\in \left( a,b\right)$, tal que
\begin{equation}
  f\left( c\right) =\frac{1}{b-a}\int_{a}^{b}f(x)dx
\end{equation}

\end{teorema}
\textit{Demostración.} Como $f$ es continua en $[a,b]$, por el teorema de valores extremos para funciones continuas\footnote{\textcite{Stewart_2012}}, $f$ alcanza mínimo y máximo en el intervalo, es decir, existen números $d$ y $e$ en el intervalo $[a,b]$ tales que $f\left(d\right) \leq f(x)\leq f(e)$ para todo $x\in[a,b]$. Por otro lado, usando la propiedad de monotonía para integrar\footnote{Ver teorema \ref{PropIntegral}}, se tiene que
\[
f(d)(b-a)\leq \int_{a}^{b}f(x)dx\leq f(e)(b-a)
\]
Dividiendo por $b-a$, resulta
\[
f\left( d\right) \leq \frac{1}{b-a}\int_{a}^{b}f(x)dx\leq f(e)
\]
luego, el valor
\[
N=\frac{1}{b-a}\int_{a}^{b}f(x)dx
\]
es un número entre dos imágenes de la función, así que, usando el teorema del valor intermedio para funciones continuas, $f$ toma todos los valores entre dos imágenes, por lo que existe un número $c$ en el intervalo $\left( a,b\right)$, tal que
\[
f(c)=N=\frac{1}{b-a}\int_{a}^{b}f(x)dx
\]
como se quería probar.

Si $f$ es una función continua y positiva en el intervalo $\left[a,b\right] $, el teorema anterior dice que el área del rectángulo con base en el intervalo $[a,b] $ en el eje $x$ y de altura $
f\left( c\right) $ es igual al área bajo la curva $y=f(x) $ en $[a,b] .$ Gráficamente, el teorema del valor medio indica que el área bajo la curva es igual al área del rectángulo con base de longitud $b-a$ y altura $f\left( c\right)$.
\[
f\left( c\right) =\frac{1}{b-a}\int_{a}^{b}f(x)dx
\]
lo que es equivalente a
\[
f\left( c\right) \left( b-a\right) =\int_{a}^{b}f(x)dx
\]

\begin{figure}[H]
\centering
\caption{Teorema del valor medio}
\includegraphics[width=7cm]{Gráficas/TeorValorMedioInteg}
\propia
\end{figure}

\begin{ejemplo}
Halle los valores de $c$ que garantiza el teorema del valor medio para integrales con la función $f(x) =3x^{2}+2x$ en el intervalo $\left[ 0,2\right]$.

\end{ejemplo}
\emph{Solución.} Dado que $f$ es continua en $\left[ 0,2\right] $, por el teorema \ref{ValorMedio} existe $c\in \left(0,2\right) $, tal que
\begin{align*}
f\left( c\right) &=\frac{1}{2}\int_{0}^{2}\left( 3x^{2}+2x\right) dx\\
\Rightarrow 3c^{2}+2c&=\frac{1}{2}\left( 3\frac{2^{3}-0^{3}}{3}+2\frac{2^{2}-0^{2}}{2}\right)\\
3c^{2}+2c&=6
\end{align*}
es decir,
\begin{align*}
c&=\frac{-2\pm \sqrt{4-\left( -72\right) }}{6}\\
&=\frac{-2\pm 2\sqrt{19}}{6}\\
c_{1}&\approx 1.1196\\
c_{2}&\approx -1.7863
\end{align*}
así que el valor de $c\in \left( 0,2\right) $ es
\[
c=\frac{-2+2\sqrt{19}}{6}\approx 1.1196
\]

\section{Teorema fundamental del cálculo}

En esta sección se presenta una de las herramientas más
importantes del cálculo, el teorema fundamental del cálculo. En su
primera parte, nos proporciona una regla de derivación, lo que nos dice
que una antiderivada de una función continua en un intervalo $[a,b] $ es $\int_{a}^{x}f(t) dt$.

En su segunda parte, el teorema
afirma que las integrales definidas se pueden calcular por medio
de antiderivadas.

\subsection{Funciónes de área de acumulación}

Sea $f $ una función continua en $[a,b]$, la \emph{función de área de acumulación} de $f$ a partir de $t=a$ es
\[
F(x) =\int_{a}^{x}f(t)dt
\]

\begin{ejemplo} Veamos los dos ejemplos siguientes:

\begin{parrafonumerado}

\item Calcule la función de área de acumulación de $f(x)=12x^{3}+4x+8$ a partir de $t=0$.

\item Halle la derivada respecto a $x$ de la función área de acumulación.
\end{parrafonumerado}

\end{ejemplo}

\emph{Solución}

\begin{parrafonumerado}
\item Usando las propiedades de la integral definida y las fórmulas para las integrales de las potencias, se tiene
\begin{align*}
  F(x) &=\int_{0}^{x}(12t^{3}+4t+8)dt\\
  &=12\frac{x^{4}-0^{4}}{4}+4\frac{x^{2}-0^{2}}{2}+8\left( x-0\right)\\
  &=3x^{4}+2x^{2}+8x
\end{align*}
que es la función de área de acumulación.

\item Como $\int_{0}^{x}f(t)dt=3x^{4}+2x^{2}+8x$, al derivar resulta
\begin{align*}
  F'(x) &=\frac{d}{dx}\int_{0}^{x}f(t)dt\\
  &=\frac{d}{dx}\left(3x^{4}+2x^{2}+8x\right) \\
  &=12x^{3}+4x+8
\end{align*}
es decir,
\[
  \frac{d}{dx} \int_{0}^{x} f(t)\,dt = f(x)
\]
\end{parrafonumerado}

\clearpage

\begin{ejemplo}\leavevmode
\begin{parrafonumerado}
\item Calcule la función de área de acumulación de $g(x) =e^{x}+2^{x}$ a partir de $t=0$.

\item Halle la derivada respecto a $x$ de la funcíon de área de acumulación.
\end{parrafonumerado}
\end{ejemplo}

\emph{Solución}

\begin{parrafonumerado}
\item Similarmente al ejemplo anterior,
\begin{align*}
  G(x)
  &=\int_{0}^{x}(e^{t}+2^{t})dt\\
  &=\int_{0}^{x}e^{t}dt+\int_{0}^{x}2^{t}dt\\
  &=e^{x}-e^{0}+\frac{2^{x}-1}{\ln 2}\\
  &=e^{x}+\frac{2^{x}}{\ln 2}-1-\frac{1}{\ln 2}\\[1em]
  G'(x) &=\frac{d}{dx}\int_{0}^{x}g(t)dt\\
  &=\frac{d}{dx}\left(e^{x}+\frac{2^{x}}{\ln 2}-1-\frac{1}{\ln 2}\right)\\
  &=e^{x}+\frac{2^{x}}{\ln 2}\ln 2\\
  &=e^{x}+2^{x}\\
  &=g(x)
\end{align*}

\end{parrafonumerado}

El TFC, en su parte I, es la generalización de los dos ejemplos anteriores. Nos dice que, si formamos una función de $x$ por medio de la integral de $f$ en el intervalo $[a,x]$, entonces la derivada con respecto a $x$ de esa función integral es igual a $f$ evaluada en $x$. O también que $\int_{a}^{x}f(t)dt$ es una antiderivada de la función $f(x)$.

\subsection{Teorema fundamental del cálculo, parte I}

Sea $f$ una función continua en el intervalo $[a,b] $ y $F(x)=\int_{a}^{x}f(t)dt$. Entonces
\[
F'(x)=\frac{d}{dx}\left(\int_{a}^{x}f(t)dt\right)=f(x)
\]
para todo $x\in \left( a,b\right)$.

\textit{Demostración}. Notemos que, por el teorema \ref{PropIntegral},
\[
F(x+h)-F(x)=\int_{a}^{x+h}f(t)dt-\int_{a}^{x}f(t)dt=\int_{x}^{x+h}f(t)dt
\]

Ahora bien, dado que $f$ es continua en $[a,b]$, entonces también es continua en $\left[ x,x+h\right]$ para $a\leq x,\ x+h\leq b$. Por otro lado, por el teorema de los valores extremos\footnote{\textcite{Zill_Wright_2011}}, existen $c,d\in \lbrack x,x+h]$ tales que para todo $t$ en este intervalo $f\left( c\right) \leq f(t)\leq f(d)$, y calculando la integral en el intervalo $[x,x+h],$ resulta que $f(c)h\leq \int_{x}^{x+h}f(t)dt\leq f(d)h,$ o también al dividir entre $h$ (suponemos que $h>0$):
\begin{equation}\label{TeorFCParteI}
f(c)\leq \frac{1}{h}\int_{x}^{x+h}f(t)dt\leq f(d)
\end{equation}

\enlargethispage*{\baselineskip}

Además, si $h\rightarrow 0$, entonces $c\to x$ y $d\to x$, es decir,
\[
\lim_{h\rightarrow 0}f(c)\leq \lim_{h\rightarrow 0}\frac{1}{h}
\int_{x}^{x+h}f(t)dt\leq \lim_{h\rightarrow 0}f(d)
\]
y
\begin{equation}\label{TeorFCPIlimites}
f(x) \leq \lim_{h\rightarrow 0}\frac{1}{h}\int_{x}^{x+h}f(t)dt\leq f(x)
\end{equation}

En este sentido, si $h<0$, entonces la desigualdad \eqref{TeorFCParteI} cambia de sentido, pero, de todas maneras, obtenemos el resultado \eqref{TeorFCPIlimites}.

Y para finalizar,
\begin{align*}
F'(x)=\lim_{h\rightarrow 0}\frac{ F(x+h)-F(x)}{h}=\lim_{h\rightarrow 0}\frac{1}{h}\int_{x}^{x+h}f(t)dt=f(x)
\end{align*}
como se quería probar.

Tengamos en cuenta que, en general, si
\[
F(x)=\int_{a}^{g(x)}f(t)dt
\]
al usar la regla de derivación que nos proporciona el TFC, parte I, combinada con la regla de la cadena
para derivadas, se obtiene que
\[
F'(x)=\frac{d}{dx}\int_{a}^{g(x)}f(t)dt=f(g(x))\frac{d(g(x))}{dx}
\]
algo que se usa en el siguiente ejemplo.

\begin{ejemplo}
Sea $F(x) =\int_{-x}^{3}\tan ^{3}zdz$. Encuentre la derivada de $F$.
\end{ejemplo}

\emph{Solución.} Por propiedades de la integral (teorema \ref{PropIntegral}),
\[
F(x) =\int_{-x}^{3}\tan ^{3}zdz=-\int_{3}^{-x}\tan
^{3}zdz
\]
y por el TFC, parte I,
\[
F'(x) =-\tan ^{3}(-x) \times \left( -1\right)
=\tan ^{3}(-x)
\]

\begin{ejemplo} Si $F(x) =\int_{3}^{x^{4}}\frac{1}{t^{8}+1}dt$, halle $F'\left(1\right)$.
\end{ejemplo}

\emph{Solución.} Por TFC,
\[
F'(x)=\frac{1}{(x^8)^4+1}\times(4x^3)=\frac{ 4x^{3}}{x^{32}+1}
\]
entonces
\[F'\left( 1\right) =\frac{4}{1^{32}+1}= 2\]

\begin{ejemplo}
Halle una función $f$ y un número $a$ tal que
\[
\int_{a}^{x}\frac{\left( f(t) \right) ^{3}}{t}dt=9x^{3}-9
\]
\end{ejemplo}

\emph{Solución.} Para hallar la función $f$, derivamos con respecto a $x$ la ecuación dada y obtenemos
\begin{gather*}
\begin{aligned}
\frac{d}{dx}\int_{a}^{x}\frac{\left( f(t) \right) ^{3}}{t}dt&=\frac{d}{dx}\left( 9x^{3}-9\right)\\
\Rightarrow \frac{\left( f(x)\right) ^{3}}{x}&=27x^{2}\\
\Rightarrow \left( f(x) \right)^{3}&=27x^{3}
\Rightarrow f(x) = 3x
\end{aligned}
\end{gather*}

Para hallar $a$, reemplazamos en la ecuación inicial $x$ por $a$, con lo que se obtiene
\begin{align*}
\int_{a}^{a}\frac{\left( f(t) \right) ^{3}}{t}dt&=9a^{3}-9\\
0&=9a^{3}-9\Rightarrow\ 9a^{3}=9\Rightarrow\ a^{3}=1
\end{align*}
Por tanto, $a=1$. Usamos acá que $\int_a^af(x)\,dx=0$.

\begin{ejemplo}
Halle una función $f$ y un número $a$, tal que $\int_{2x}^{a}\left(-\frac{f(t) }{t^{2}}\right)dt=x^{2}-81$.
\end{ejemplo}

Para escribir la integral en la forma en que la usa el TFC, inicialmente intercambiamos límites de integración:
\[
\int_{2x}^{a}\left(-\frac{f(t) }{t^{2}}\right)dt=-\int_{a}^{2x}\left(-\frac{f(t) }{t^{2}}\right)dt=\int_{a}^{2x}\frac{f(t) }{t^{2}}dt
\]

\enlargethispage*{\baselineskip}

Ahora derivamos cada miembro de la ecuación dada con respecto a $x$:
\begin{gather*}
\begin{aligned}
\frac{d}{dx}\left( \int_{a}^{2x}\frac{f(t) }{t^{2}}dt\right) &=\frac{d}{dx}\left( x^{2}-81\right) \Rightarrow\ \frac{f\left( 2x\right)}{\left( 2x\right) ^{2}}\times 2=2x \Rightarrow\ f\left( 2x\right) =4x^{3}
\end{aligned}
\end{gather*}

Sustituyendo $u=2x$, es decir, $x=\frac{u}{2}$, tenemos
\[
f\left(u\right)=4\left(\frac{u}{2}\right)^{3}=\frac{4u^{3}}{8}=\frac{u^{3}}{2}
\]
es decir, la función buscada es $f(u)=\frac{u^3}{2}$.

Por otro lado, si reemplazamos en la ecuación inicial $x$ por $\cfrac{a}{2}$, tenemos
\[
\int_{a}^{a}\left(-\frac{f(t) }{t^{2}}\right)dt=\left(\frac{a}{2}\right)^{2}-81\Rightarrow 0=\frac{a^2}{4}-81\Rightarrow a=\pm 18
\]
y por tanto, un número puede ser $a=18$.

\subsection{Teorema fundamental del cálculo, parte II}

Si $f$ es una función continua en $[a,b]$, entonces
\[
\int_{a}^{b}f(x)\, dx=F(b)-F(a)
\]
para $F$ una antiderivada de $f$.

\textit{Demostración.} Sea $F$ una antiderivada cualquiera de $f$. Por el TFC en su parte I, $A(x)=\int_{a}^{x}f(t)dt$ es otra antiderivada de $f$; luego, como la diferencia de dos antiderivadas es una constante, se tiene que $F(x)-A(x)=C$; entonces, $F(x)=A(x) +C$; por tanto,
\begin{align*}
F(b)-F(a)&=\left( \int_{a}^{b}f(t)dt+C\right) -\left(\int_{a}^{a}f(t)dt+C\right) \\
&=\int_{a}^{b}f(t)dt
\end{align*}
como se quería probar.

\clearpage

\begin{ejemplo}
Cálculo de integrales definidas:
\end{ejemplo}

\begin{parrafonumerado}
\item $\displaystyle{\int_{0}^{\frac{\pi }{4}}(\sen (2x)+\cos (3x))dx}$

\vspace{\baselineskip}

\emph{Solución.} Una antiderivada de $\sen (2x)+\cos (3x)$ es
\begin{align*}
F(x) & =\int (\sen (2x)+\cos (3x))dx\\
& =-\frac{1}{2}\cos 2x+\frac{1}{3}\sen (3x)
\end{align*}

y por el TFC, parte I,
\begin{align*}
\int_{0}^{\frac{\pi }{4}}&(\sen (2x)+\cos (3x))\,dx=F(\pi/4)-F(0)\\
&=-\frac{1}{2}\cos(\pi/2)+\frac{1}{3}\sen (3\pi/4)-\left(-\frac{1}{2}\cos 0+\frac{1}{3}\sen 0\right)\\
&=\frac{1}{3}\frac{\sqrt{2}}{2}+\frac{1}{2}\\
&=\frac{\sqrt{2}}{6}+\frac{1}{2}
\end{align*}

\item $\displaystyle{\int_{0}^{\frac{\pi }{8}}\sec (2x)\tan (2x)dx}$

\vspace{\baselineskip}

\emph{Solución}
\begin{align*}
\int_{0}^{\frac{\pi }{8}}\sec (2x)\tan (2x)\,dx&=\int \sec(2x)\tan(2x)\,dx \Big|_{0}^{\frac{\pi }{8}}\\
&=\frac{1}{2}\sec(2x)\Big|_{0}^{\frac{\pi }{8}}\\
&=\frac{1}{2}\sec \frac{\pi }{4}-\frac{1}{2}\sec0\\
&=\frac{\sqrt{2}}{2}-\frac{1}{2}
\end{align*}

\item $\displaystyle{\int_{\frac{\pi }{12}}^{\frac{\pi }{9}}(\sec ^{2}(3x)+\csc^{2}(3x))dx}$

\vspace{\baselineskip}

\emph{Solución}
\begin{gather*}
\begin{aligned}
  \int_{\frac{\pi }{12}}^{\frac{\pi }{9}}\bigl(\sec ^{2}(3x)&+\csc^{2}(3x)\bigr)dx\\
  &=\left(\int \sec ^{2}(3x)dx+\int \csc ^{2}(3x)dx\right)\Bigg|_{\pi/12}^{\pi/9}\\
  &=\left(\frac{1}{3}\tan 3x-\frac{1}{3}\cot 3x\right)\Bigg|_{\pi/12}^{\pi/9}\\
  &=\frac{1}{3}\left(\underset{x=\frac{\pi}{9}}
    {\underbrace{\left(\tan(\pi/3)-\cot(\pi/3)\right)}}
  -\underset{x=\frac{\pi }{12}}{\underbrace{(\tan(\pi/4)-\cot(\pi/4))}}\right)\\
  &=\frac{2}{9}\sqrt{3}
\end{aligned}
\end{gather*}

\item $\displaystyle{\int_{0}^{1}\left(\sqrt{x}+\sqrt[5]{x^{2}}+\sqrt[7]{x^{8}}\right)\,dx}$

\vspace{\baselineskip}

\emph{Solución}
\begin{gather*}
\begin{aligned}
  \int_{0}^{1}\left( \sqrt{x}+\sqrt[5]{x^{2}}+\sqrt[7]{x^{8}} \right) \,dx&=\int_{0}^{1} \left( x^{\frac{1}{2}}+x^{\frac{2}{5}}+x^{\frac{8}{7}}\right)dx \\
  &=\left(\frac{1}{\frac{3}{2}}x^{\frac{3}{2}}+\frac{1}{\frac{7}{5}}x^{\frac{7}{5}}+\frac{1}{\frac{15}{7}}x^{\frac{15}{7}}\right)\bigg|_{0}^{1}\\
  &=\underset{x=1}{\underbrace{\frac{3}{2}+\frac{5}{7}+\frac{7}{15}}}-\underset{x=0}{\underbrace{0}}\\[0.2em]
  &=\frac{70+75+49}{105}\\
  &=\frac{194}{105}
\end{aligned}
\end{gather*}

\item $\displaystyle{\int_{0}^{\frac{\pi }{4}}\frac{1+\sen x}{1-\sen x}dx}$

\vspace{\baselineskip}

\emph{Solución.} Primero transformamos la función por integrar mediante identidades:
\begin{gather*}
\begin{aligned}
  \frac{1+\sen x}{1-\sen x}&=\frac{1+\sen x}{1-\sen x}\cdot \left(\frac{1+\sen x}{1+\sen x}\right)=\frac{1+\sen x+\sen x+\sen ^{2}x}{1+\sen x-\sen x-\sen ^{2}x}\\
  &=\frac{1+2\sen x+\sen ^{2}x}{1-\sen ^{2}x}=\frac{1+2\sen x+1-\cos ^{2}x}{\cos ^{2}x}\\
  &=\frac{2}{\cos ^{2}x}+\frac{\sen x}{\cos ^{2}x}-\frac{\cos ^{2}x}{\cos ^{2}x}=2\sec ^{2}x+2\cdot \frac{1}{\cos x}\cdot \frac{\sen x}{\cos x}-1\\
  &=2\sec^{2}x+2\sec x \tan x-1
\end{aligned}
\end{gather*}

Ahora calculemos la integral:
\begin{align*}
\int_{0}^{\frac{\pi }{4}}\frac{1+\sen x}{1-\sen x}dx&= \int_{0}^{\pi/4} \left(
2\sec ^{2}x+2\sec x\tan x-1\right) dx\\
&=\left( 2\tan x+2\sec x-x\right)\big|_{0}^{\pi/4}\\
&=\underset{x=\frac{\pi }{4}}{\underbrace{2\tan \frac{\pi }{4}+2\sec \frac{\pi }{4}-\frac{\pi }{4}}}\underset{x=0}{-\underbrace{(0+2-0)}}\\
&=2+2\sqrt{2}-\frac{\pi }{4}-2=2\sqrt{2}-\frac{\pi }{4}
\end{align*}

\end{parrafonumerado}

\begin{ejemplo}
Use el TFC para derivar cada una de las siguientes funciones:
\end{ejemplo}

%%\begin{multicols}{4}
\begin{parrafonumerado}

\item $g(x)=\displaystyle{\int_{-1}^{x}\frac{2}{\sqrt{t^{2}+t+1}}dt}$

\item $g(y)=\displaystyle{\int_{y}^{2}\frac{\sen x}{\sqrt{x^{2}+1}}dx}$

\item $g(z)=\displaystyle{\int_{3}^{z^{2}+2z}\sqrt[3]{x^{4}+x}dx}$

\item $g(s)=\displaystyle{\int_{-\cos s}^{\tan s}\frac{t^{3}}{\sqrt{t^{2}+1}}dt}$
\end{parrafonumerado}
%%\end{multicols}

\emph{Solución}

\begin{parrafonumerado}

\item $\displaystyle{g'(x)=\frac{d}{dx}\int_{-1}^{x}\frac{2}{\sqrt{t^{2}+t+1}}dt
=\frac{2}{\sqrt{x^{2}+x+1}}}$

\item $\displaystyle{g^\prime(y)
=\frac{d}{dy}\int_{y}^{2}\frac{\sen x}{\sqrt{x^{2}+1}}dx
=-\frac{d}{dy}\int_{2}^{y}\frac{\sen x}{\sqrt{x^{2}+1}}dx}
=-\frac{\sen y}{\sqrt{y^{2}+1}}$

\item $\displaystyle{g'(z)=\frac{d}{dz}\int_{3}^{z^{2}+2z}\sqrt[3]{x^{4}+x}dx=\sqrt[3]{%
\left( z^{2}+2z\right) ^{4}+z^{2}+2z}\cdot \left( 2z+2\right)}$

\item Inicialmente, reescribimos la integral de la siguiente manera:
\begin{align*}
g(s)&=\int_{-\cos s}^{\tan s}\frac{t^{3}}{\sqrt{t^{2}+1}}dt=\int_{-\cos
s}^{0}\frac{t^{3}}{\sqrt{t^{2}+1}}dt+\int_{0}^{\tan s}\frac{t^{3}}{\sqrt{%
t^{2}+1}}dt\\
&=\int_{0}^{\tan s}\frac{t^{3}}{\sqrt{t^{2}+1}}
dt-\int_{0}^{-\cos s}\frac{t^{3}}{\sqrt{t^{2}+1}}dt
\end{align*}
Entonces,
\begin{align*}
g'(s)&=\frac{d}{ds}\int_{0}^{\tan s}\frac{t^{3}}{\sqrt{t^{2}+1}}dt-\frac{d}{ds}\int_{0}^{-\cos s}\frac{t^{3}}{\sqrt{t^{2}+1}}dt\\
&=\frac{\tan ^{3}s}{\sqrt{\tan ^{2}s+1}}\cdot \sec ^{2}s-\frac{-\cos ^{3}s}{\sqrt{\cos ^{2}s+1}}\cdot \sen s\\
&=\frac{\tan ^{3}s\sec ^{2}s}{\sqrt{\tan ^{2}s+1}}+\frac{\cos^{3}s\cdot \sen s}{\sqrt{\cos ^{2}s+1}}
\end{align*}
\end{parrafonumerado}

\section{Ejercicios del capítulo~\thechapter}

\group

En los ejercicios \ref{AntiInici}-\ref{AntiFin}, encuentre el valor de las integrales:

\exercise{\label{AntiInici}$\displaystyle\int (5-6x)dx$}{$5x-3x^{2}+C$}

\exercise{$\displaystyle\int \left( \frac{t^{2}}{2}+4t^{3}\right)dt$}{$\frac{t^{3}}{6}+t^{4}+C$}

\exercise{$\displaystyle\int \left( \frac{1}{5}-\frac{2}{x^{3}}+2x\right) dx$}{$\frac{1}{5x^{2}}\left( 5x^{4}+x^{3}+5\right) +C$}

\exercise{$\displaystyle\int \left( \frac{\sqrt{x}}{2}+\frac{2}{\sqrt{x}}\right) dx$}{$4\sqrt{x}+\frac{1}{3}x^{\frac{3}{2}}+C$}

\exercise{$\displaystyle\int \left( \frac{1}{7}-\frac{1}{y^{\frac{5}{4}}}\right) dy$}{$\frac{\sqrt[4]{y^{5}}+28}{7\sqrt[4]{y}}+C$}

\exercise{$\displaystyle\int \frac{4+\sqrt{t}}{t^{3}}dt$}{$-\frac{2}{3t^{2}}\left( \sqrt{t}+3\right) +C$}

\exercise{$\displaystyle\int \frac{5a^{x}dx}{6}$}{$\frac{5a^{x}}{6\ln (a)}+C$}

\exercise{$\displaystyle\int \left( \frac{-\sec ^{2}x}{3}\right)dx$}{$-\frac{\tan x}{3}+C$}

\exercise{$\displaystyle\int \frac{2}{5}\sec (\theta )\tan (\theta)d\theta $}{$\frac{2}{5}\sec (\theta )+C$}

\exercise{$\displaystyle\int \left( 5x^{7}+6x^{4}+7x-3\right) dx$}{$\frac{5}{8}x^{8}+\frac{6}{5}x^{5}+\frac{7}{2}x^{2}-3x+C$}

\exercise{$\displaystyle\int\left( 6e^{2x}+2e^{\frac{1}{3}x}\right) dx$}{$3e^{2x}+6e^{\frac{1}{3}x}+C$}

\exercise{$\displaystyle\int \left(2^{x}+3^{x}+4^{2x}+5x^{-1}\right) dx$}{$\frac{2^{x}}{\ln 2}+\frac{3^{x}}{  \ln 3}+\frac{16^{x}}{\ln 16}+5\ln \left\vert x\right\vert +C$}

\exercise{$\displaystyle\int \left( \sen x+\cos x\right) ^{2}dx$}{$x-\frac{1}{2}\cos 2x+C$}

\exercise{$\displaystyle\int \left( \sec x-\cos x\right) ^{2}dx$}{$\tan x-\frac{3}{2}x+\frac{1}{4}\sen (2x)+C$}

\exercise{$\displaystyle\int \left(\csc x+\sen x\right) ^{2}dx$}{$-\cot \left( x\right) +\frac{5x}{2}-\frac{\sin \left(2x\right)}{4}+C$
}

\exercise{$\displaystyle\int \left( \tan x+\cot x\right) ^{2}dx$}{$\tan x-\cot x+C$}

\exercise{$\displaystyle\int \frac{1+\sen x}{1-\sen x}dx$}{$-\cot \left( 2x\right) -\csc 2x-x+C$}

\exercise{$\displaystyle\int \left( 2+\tan ^{2}\theta \right)d\theta $}{$\theta +\tan \theta +C$}

\exercise{$\displaystyle\int \left( 1-\cot ^{2}(x)\right) dx$}{$2x+\cot (x)+C$}

\exercise{$\displaystyle\int \frac{\csc (\theta )}{\csc (\theta)-\sen (\theta )}d\theta $}{$\tan \left( \theta \right) +C$}

\exercise{$\displaystyle\int x^{-3}(x+1)dx$}{$-\frac{1}{x^{2}}\left( x+\frac{1}{2}\right) +C$}

\exercise{$\displaystyle\int \left(x^{-3}+4\sqrt[3]{x}-\frac{7}{x^{5}}+4\sqrt{x}\right)dx$}
{$-\frac{1}{2x^{2}}+3\sqrt[3]{x^{4}}+\frac{7}{4x^{4}}+\frac{8}{3}\sqrt{x^{2}}+C$}

\exercise{$\displaystyle\int \left( \sen x+\cos x+\sen \left(3x\right) -\cos \left( \frac{1}{4}x\right) \right) dx$}{$-\cos x+\sen x-
\frac{1}{3}\cos \left( 3x\right) -4\sen \left( \frac{1}{4}x\right) +C$}

\exercise{$\displaystyle\int \left( \sec ^{2}x+4\csc ^{2}\left(2x\right) +8\sec ^{2}\left( 4x\right) \right) dx$}{$\tan x-2\cot \left(
2x\right) +2\tan \left( 4x\right) +C$}

\exercise{$\displaystyle\int \left( \sec x\tan x-4\csc 2x\cot2x\right) dx$}{$\sec x+2\csc \left( 2x\right) +C$}

\exercise{\label{AntiFin}$\displaystyle\int \frac{1}{2}\left( \csc ^{2}(\theta)-\csc (\theta )\cot (\theta )\right) d\theta $}{$\frac{-\cot (\theta)+\csc (\theta )}{2}+C=F(\theta )$}

%%%%%%%%%%%%%%%%%%%%%%%%%%%%%%%%%%%%%%

\vspace{\parskip}

\exercise{Sea $v(t)$ la función de velocidad de una partícula en un instante $t$ en segundos.  Hallar la función de posición $s(t)$ de la partícula en cada caso}

\subexercise{$v(t)=t^{3}+t^{2}+t,\  s\left(0\right)=3$}{$s\left( t\right) =\frac{t^{4}}{4}+\frac{t^{3}}{3}+\frac{t^{2}}{2}+3$}

\subexercise{$v(t)=\sen t+\cos t,\  s\left(\pi\right)=2$}{$s\left( t\right) =\cos \left( t\right) -\sen \left( t\right)
+1$}

\subexercise{$v(t)=\sec ^{2}t+\csc ^{2}t,\
s\left( \frac{\pi }{4}\right) =2$}{$s\left( t\right) =\tan \left( t\right)
-\cot \left( t\right) -\frac{\sqrt{2}}{2}+3$}

\subexercise{$v(t)=\frac{1-\cos ^{2}t}{\sen ^{2}t},\  s\left( \frac{\pi }{4}\right) =8$}{$s\left( t\right) =t+8-\frac{\pi
}{4}$}

\exercise{Sea $a(t)$ la función de aceleración de una partícula en un instante. Hallar la función de posición $s(t)$ de la partícula ($t$ en segundos $s$ en metros) en cada caso}

\subexercise{$a(t)=t^{3}+t^{2},\  v\left( 0\right)
=3,s\left( 0\right) =5$}{$s\left( t\right) =\frac{t^{5}}{20}+\frac{t^{4}}{12}+3t+5$}

\subexercise{$a(t)=\sen t+\cos t,\  v(\pi )=4,$ $\
s\left( \pi \right) =2$}{$s\left( t\right) =-\sen \left( t\right) -\cos
\left( t\right) +3t-3\pi +1$}

\subexercise{$a(t)=\frac{2t^{3}+t^{2}}{t^{5}},\
v\left( 1\right) =4,s\left( 1\right) =10$}{$s\left( t\right) =-2\ln\left\vert t\right\vert+\frac{1}{2t}+\frac{13}{4}t^{2}+\frac{25}{4}$}

\subexercise{$a(t)=-9.8t,\  v(2)=7,$ $\ s\left(3\right) =9\,$}{$s\left( t\right) =-\frac{9.8}{6}t^{3}+26.6t-26.7$}

\exercise{En cada caso encuentre la función que satisfaga las condiciones dadas}

\subexercise{Pasa por el punto $(3,4)$ y en el punto $(x,f(x))$ la pendiente de la recta tangente es$\ 3x^{2}-8$.}{$f(x)=x^{3}-8x+19$}

\subexercise{Pasa por el punto $(0,3)$ y en el punto $(x,f(x))$ la pendiente de la recta tangente es $e^{x}-x$.}{$f(x)=e^{x}-
\frac{x^{2}}{2}+2$}

\exercise{Encuentre la función que cumple:}

\subexercise{$f^{\prime\prime}\left( x\right) =e^{x}-x,\ f\left( 1\right) =3,\ f^{\prime}\left( 2\right) =1$.}{$f(x)=e^{x}-\frac{x^{3}}{6}
+(3-e^{2})x+\frac{6e-1-6e^{2}}{6}$}

\subexercise{$f^{\prime\prime\prime}\left( x\right) =\sen x-\cos\left( 2x\right),\  f\left( 0\right) =f^{\prime}\left( \frac{\pi }{2}
\right) =f^{\prime\prime}\left( \pi \right) =1$.}{$f\left( x\right) =-\cos x-\frac{\sen \left( 2x\right) }{8}+x^{2}+\frac{7}{4}x+2$}

\subexercise{$\cfrac{d^{\,2}y}{dx^{2}}=2x^{2}+3x^{3},\ y(1)=3,y^{\prime}\left(2\right)=3$.}{$y=\frac{x^{4}}{12}+\frac{3x^{5}}{20}-\frac{43}{3}x+\frac{156}{5}$}

\subexercise{$\cfrac{d^{\,4}y}{dx^{4}}=120x^{2},\ y\left( 0\right) =y^{\prime}\left( 0\right) =y^{\prime\prime\prime}\left( 1\right) =1$.}{$
y\left( x\right) =8\cos \left( \frac{1}{2}x\right) -\sen x+\frac{1}{8}e^{2x}+x^{4}+\frac{11}{4}x^{2}+\frac{15}{4}x-\frac{41}{8}$}

\exercise{La velocidad de una partícula es $v(t)=\frac{1}{1+\cos \left( 2t\right)}$, $t$ en metros, $s$ en segundos. Si
su posición en $t=\frac{\pi }{8}$ es $2$ metros, hallar la función de posición de la partícula.}{La función de posición de la partícula es: $s(t)=-\cot \left( 2t\right) +\frac{1}{2}\csc \left(2t\right) +3-\frac{\sqrt{2}}{2}$}

\exercise{La aceleración de una partícula es $a(t)=\sen \left( 2t\right) +e^{2t}+2^{4t}$, con $t$ en metros, $s$ en
segundos, su velocidad inicial es 6 metros por segundo y su posición inicial es 4 metros. Hallar la función de posición de la partícula.}{$s\left( t\right) =-\frac{\sen \left( 2t\right) }{4}+\frac{e^{2t}}{4}+\frac{2^{4t}}{16\ln ^{2}\left( 2\right) }+\left( 6-\frac{1}{4\ln \left(2\right) }\right) t+\frac{15}{4}-\frac{1}{16\ln ^{2}\left( 2\right) }$}

\exercise{La pendiente de la tangente a la gráfica de $f$ en $\left(x,f\left( x\right) \right)$ es $\frac{\cos  ^{2}x}{1+\sen x}$ y $f\left( \frac{\pi }{6}\right) =4$. Hallar la función $f\left( x\right)$.}{$f\left( x\right) =x+\cos x+4-\frac{\pi}{6}-\frac{\sqrt{3}}{2}$}

\vspace*{\parskip}

En los siguientes ejercicios para cada una de las funciones $f(x)$ dadas, hallar el valor de $c$ que garantiza el teorema del valor medio para integrales en los intervalos indicados y hacer una gráfica de la situación.

\begin{multicols}{2}
\exercise{$f\left( x\right) =6x^{2}+4x+4$ en $\left[1,3\right].$}{$c=\frac{2}{3}\sqrt{13}-\frac{1}{3}$}

\exercise{$f\left( x\right) =x^{3}$ en $\left[ 2,4\right] .$}{$c=\sqrt[3]{30}=3.1072$}

\exercise{$f\left( x\right) =3^{x}$ en $\left[ 0,2\right] $.}{$c=\frac{\ln \left( \frac{4}{\ln 3}\right) }{\ln 3}=1.1763$}

\exercise{$f\left( x\right) =2^{3x}$ en $\left[ 0,4\right] $.}{$c=\frac{\ln \left( \frac{1365}{4\ln 2}\right) }{\ln 8}=2.9812$}

\exercise{$f\left( x\right) =2\sen x$ en $\left[ 0,\frac{\pi }{2}\right] $.}{$c=\arcsin \left( \frac{2}{\pi }\right) =0.69011$}

\exercise{$f\left( x\right) =4\cos x$ en $\left[ 0,\frac{\pi }{2}\right] $.}{$c=\arccos \left( \frac{2}{\pi }\right)=0.880\,69$}
\end{multicols}

\vspace*{\parskip}

Para cada una de las funciones $f\left( x\right) $ dadas en los siguientes ejercicios: a) Calcular la función $A\left( x\right) =\int_{a}^{x}f\left( t\right) dt$, de área de acumulación de $f\left( x\right) $ a partir del valor de $t=a$ indicado. b) Hallar la derivada respecto a $x$ de la función área de acumulación $\frac{dA}{dx}$.

\exercise{$f\left( x\right) =3x^{2}+6x+2$ a partir de $t=2$}{$a)$ $A\left( x\right) =x^{3}+3x^{2}+2x-24\qquad b)$ $\frac{dA}{dx}
=3x^{2}+6x+2=f\left( x\right) $}

\exercise{$f\left( x\right) =4x^{3}+3x^{2}+1$ a
partir de $t=3$}{$a)$ $A\left( x\right) =x^{4}+x^{3}+x-111\qquad b)$ $
\frac{dA}{dx}=4x^{3}+3x^{2}+1=f\left( x\right) $}

\exercise{$f\left( x\right) =12x^{5}+6x^{2}+2$ a partir de $t=1$}{$a)$ $A\left( x\right) =2x^{6}+2x^{3}+2x-6\qquad b)$ $ \frac{dA}{dx}=12x^{5}+6x^{2}+2=f\left( x\right) $}

\exercise{$f\left( x\right) =e^{x}+2x$ a partir de $t=0.$}{$a)$ $A\left( x\right) =e^{x}+x^{2}-1\qquad b)$ $\frac{dA}{dx}=e^{x}+2x=f\left( x\right) $}

\exercise{$f\left( x\right) =2^{x}+3^{x}+8x+1$ a
partir de $t=0.$}{$a)$ $A\left( x\right) =\frac{2^{x}}{\ln 2}+\frac{3^{x}}{%
\ln 3}+4x^{2}+x-\frac{1}{\ln 2}-\frac{1}{\ln 3}\qquad b)$ $\frac{dA}{dx}%
=2^{x}+3^{x}+8x+1=f\left( x\right) $}

\exercise{$f\left( x\right) =2\sen x+3x$ a partir de $t=0.$}{$a)$ $A\left( x\right) =\frac{3}{2}x^{2}-2\cos x+2\qquad b)$ $\frac{%
dA}{dx}=2\sen x+3x=f\left( x\right) $}

\exercise{$f\left( x\right) =e^{2x}+\sen x+1$ a
partir de $t=0.$}{$a)$ $A\left( x\right) =x-\cos x+\frac{1}{2}e^{2x}+\frac{%
1}{2}\qquad b)$ $e^{2x}+\sen x+1=f\left( x\right) $}

\exercise{$f\left( x\right) =e^{4x}+\sen x+\cos x$ a
partir de $t=0.$}{$a)$ $A\left( x\right) =\sen x-\cos x+\frac{1}{4}e^{4x}+\frac{3}{4}\qquad b)$ $\frac{dA}{dx}=e^{4x}+\sen x+\cos x=f\left( x\right) $}

\exercise{Encuentre la derivada de cada una de las siguientes funciones}
\begin{multicols}{2}
\subexercise{$f(x)=\displaystyle\int_{-3}^{x^{4}}\sqrt{t}dt$}{$4x^{3}\sqrt{x^{4}}$}

\subexercise{$g(y)=\displaystyle\int_{\cos y}^{\sin y}\sqrt{t^{2}+1}dt$}{$\sqrt{\cos
^{2}y+1}(\sin y)+\sqrt{\sin ^{2}y+1}(\cos y)$}

\subexercise{$h(z)=\displaystyle\int_{z^{2}}^{z^{4}}\frac{\sqrt{t}}{t}dt$}{$\frac{%
-2\sqrt{z^{2}}+4\sqrt{z^{4}}}{z}$}

\subexercise{$p(x)=\displaystyle\int_{-x}^{\cos (x^{4})}\frac{\sqrt{t}}{t+1}dt$}{$\frac{\sqrt{-x}}{-x+1}-\frac{%
4x^{3}\sin (x^{4})\sqrt{\cos (x^{4})}}{\cos (x^{4})+1}$}
\end{multicols}

\exercise {En cada caso hallar una función $f$ y un número $a$ que cumplan la igualdad dada}

\subexercise{$\displaystyle\int_{a}^{x}\frac{f\left( t\right) }{t^{2}}dt=3x^{2}-9$}{$f(x)=6x^{3}\qquad a=\pm \sqrt{3}$}

\subexercise{$\displaystyle\int_{a}^{x}\frac{f(t)+1}{f(t)-1}dt=\frac{x-1}{2x+1}$}{\linebreak$f(x)=-3-(2x+1)^{2}\qquad  a=1$}

\subexercise{$\displaystyle\int_{2x}^{a-1}\left( t+1\right) ^{2}f(t)dt=\frac{x^{3}}{2x+1}$
}{$f(z)=\frac{2z^{3}+3z^{2}}{8(z+1)^{4}}\qquad a=-1$}

\exercise {Calcular las siguientes integrales definidas}

\subexercise{$\displaystyle\int_{-2}^{2}\left( 4x^{3}+6x^{2}+8x-1\right) dx$} {$28$}

\subexercise{$\displaystyle\int_{0}^{\pi }\sin \left( \frac{1}{4}x\right) dx$=}{$4-2\sqrt{2}$}

\subexercise{$\displaystyle\int_{0}^{\pi }\left( \cos \left( \frac{1}{2}x\right) +\sin \left( \frac{1%
}{2}x\right) \right) ^{2}dx$}{$\pi +2$}

\subexercise{$\displaystyle\int_{0}^{1}\left( \sqrt[3]{8t}-5t^{\frac{3}{4}}+7\sqrt[7]{t}\right) dt$}{$\frac{267}{56}$}

\subexercise{$\displaystyle\int_{0}^{\frac{\pi }{8}}\frac{\sin (2x)}{\cos ^{2}(2x)}dx$}{ $\frac{1}{2}  \sqrt{2}-\frac{1}{2}$}

\subexercise{$\displaystyle\int_{\frac{\pi }{2}}^{\frac{3\pi }{2}}\frac{\cos \left( \frac{x}{2}%
\right) }{\sin ^{2}\left( \frac{x}{2}\right) }dx$}{ $0$}

\subexercise{$\displaystyle\int_{0}^{\frac{1}{2}}\frac{1+\sqrt{1-x^{2}}}{\sqrt{1-x^{2}}}dx\qquad $ }{$\frac{1}{6}\pi +\frac{1}{2}$}

\subexercise{$\displaystyle\int_{0}^{1}\frac{x^{2}+4}{x^{2}+16}dx\qquad $}{$3\arctan 4-\frac{3}{2}\pi +1$}

\subexercise{$\displaystyle\int_{-1}^{1}\frac{x^{4}+4x^{2}+2}{x^{2}+4}dx\qquad $}{ $\pi -2\arctan 2+%
\frac{2}{3}$}

\subexercise{$\displaystyle\int_{4}^{8}\frac{y+\sqrt{y^{2}-4}}{y^{2}\sqrt{y^{2}-4}}%
dy$}{$=\frac{\arcsec(4)}{2}+\frac{1}{8}-\frac{\arcsec(2)}{2}$}

\subexercise{$\displaystyle\int_{\frac{\pi }{16}}^{\frac{\pi }{8}}\left( \tan (2t)+\cot (2t)\right)
^{2}\ dt$}{$1-\frac{\tan (\frac{\pi }{8})}{2}-\frac{\cot (\frac{\pi }{8})}{2}$}

\subexercise{$\displaystyle\int_{0}^{1}\left( 2t+3\right) ^{4}dt$}{$\frac{1441}{5}$}

\exercise {Hallar el área de la región limitada por las gráficas de:}

\subexercise {$y=x^{4}+1,$ el eje $x$, el eje $y$ y la recta $x=3$}{$\frac{258}{5}u^{2}$}

\subexercise {$y=\cos(2x)$, el eje $x$, el eje $y$ y la recta $x=\frac{\pi }{8}$}{$\frac{1}{4}\sqrt{2}u^{2}$}

\subexercise {$y=\frac{x^{4}+x+1}{x^{6}},$ el eje $x$, y las rectas $x=1$ y $x=2$}{$\frac{297}{320}u^{2}$}

\subexercise {$y=\frac{x^{4}+1}{x^{2}+1},$ el eje $x$, el eje $y$ y la recta $x=1$}{$\left( \frac{1}{2}\pi - \frac{2}{3}\right) u^{2}$}

\subexercise {$ y=\frac{\sen^{2}x+1}{\cos ^{2}x},$ el eje $x$, el eje $y$ y la recta $x=\frac{\pi }{4}$}{$\left(2-\frac{1}{4}\pi \right) u^{2}$}

\subexercise {$y=\frac{1}{\sqrt{9-x^{2}}},$ el eje $x$ y la recta $x=1$}{$\frac{1}{2}\pi u^{2}$}
